\documentclass[10pt,a4paper]{amsart}
\usepackage[margin=19mm]{geometry}
\usepackage{amsmath,amssymb,mathtools}
\usepackage[scr=rsfs,cal=euler]{mathalfa}
\usepackage{xfrac}
\usepackage[unicode,hidelinks,bookmarks=false]{hyperref}
\newcommand{\ie}{i.e.\ }
\newcommand{\cf}{cf.\ }
\newcommand{\resp}{respectively\ }
\newcommand{\Subsection}{\subsection}
\providecommand{\nobreakdash}{\nobreak\hbox{-}}
\setlength{\emergencystretch}{2em}
\multlinegap0pt
\usepackage{tikz}
\usepackage{amssymb}
\usepackage{dsfont}
\usepackage{float}
\usepackage{mathtools}
\usepackage{enumerate}
\newtheorem{proposition}{Proposition}
	\def\RR{\mathbb R}
	\newcommand{\abs}[1]{\left|#1\right|}
\title{Consensus-based optimization with $\alpha$-stable jump processes}
\author{Pedro Aceves-Sanchez, Giacomo Albi, Federica Ferrarese, Michael Herty}
\date{Source: arXiv:2604.05626v1 (CC BY 4.0)}
\begin{document}
\maketitle

\noindent\emph{Source and licence.} Consensus-based optimization with $\alpha$-stable jump processes. Version arXiv:2604.05626v1 (CC BY 4.0), distributed by arXiv under the Creative Commons Attribution 4.0 International licence, http://creativecommons.org/licenses/by/4.0/, as recorded in the arXiv licence metadata for this exact version and shown on the abstract page https://arxiv.org/abs/2604.05626v1. Full licence notice: \texttt{licenses/arxiv-2604.05626-CC-BY-4.0.txt}.\par\smallskip

\noindent\textit{Editorial scope.} Counted unit: the article's proposition prop:mass\_bounds (source0.tex lines 542--561). The article's complete original proof of it is delivered with it (source0.tex lines 562--752, one \texttt{proof} environment of 191 lines). No mathematics is written, repaired or re-derived: the statement and the proof are the article's own lines, in the article's own notation, and no retained line is substituted or re-flowed.  Retained uncounted as the source-local context the proof needs: lines 329--336; lines 760--762.  Nothing was omitted from any retained span, so the package carries no omission record.  No retained line contains a citation, so the wrapper carries no bibliography and no bibliography entry.  No retained line contains a figure environment, an image-inclusion command or drawing code, so no image is delivered and the package declares no asset dependency.  Editorial formatting only, in addition: an AMS \texttt{amsart} preamble that restates the article's own declarations and macros used by the retained lines, namely the environments \texttt{proposition} on separate counters, together with the standard packages those lines need.  The retained environments print in this document as Proposition 1 in order of appearance, whereas the article prints the same items with the numbering of its own full document.  The complete native source of the article as distributed in the arXiv e-print for this exact version is bundled unchanged as \texttt{sources/197969/source0.tex}.
\begin{equation}\label{eq:fFP}
	\begin{aligned}
		\partial_t f(x,t) = 
		& - \nabla_x \cdot \big(\nu(\bar{x}_\mathcal{E}(t)-x) f(x,t)\big) + \sigma^2 \Delta\big(D^2(\bar{x}_\mathcal{E}(t),x) f(x,t)\big) \\
		& 
		\qquad - \gamma^{\alpha} (-\Delta)^{\alpha/2}\big(|D(\bar{x}_\mathcal{E}(t),x)|^\alpha f(x,t)\big),
	\end{aligned}
\end{equation}

	\begin{equation}\label{eq:bound_mass}
		\rho_t(B_r(x^\star)) \geq m_0 e^{-Pt} :=m_\star,
	\end{equation}

\begin{proposition}\label{prop:mass_bounds}
	Let $\nu,\gamma>0$. Let $f$ be a weak solution of \eqref{eq:fFP}. For $r>0$, define the mollifier
	\begin{equation}\label{eq:mollifier}
		\phi_r(x) = \begin{cases}
			\exp \left( 1- \frac{r^2}{r^2 - \vert x-x^\star\vert^2} \right)  \qquad \text{if } \quad \vert x-x^\star \vert < r,\\
			0  \qquad  \qquad \qquad \qquad \qquad  \text{elsewhere}.
		\end{cases}
	\end{equation}
	Assume there exists $B>0$ such that $
	\vert \bar{x}_\mathcal{E}(t) - x^\star \vert \leq B$,
	for any $t \in [0,T]$, $T>0$ (follows from quantitative Laplace principle). Then, for any $t\in[0,T]$ we have 
	\begin{equation}\label{eq:mass_bound}
		\rho_t(\mathcal{B}_r(x^\star)) \geq  m_0 e^{-P t}:= m_\star , 
	\end{equation}
	with $m_0 = \rho_o(\mathcal{B}(x^\star))$ and 
	\begin{equation*}
		P = \nu \frac{C_{drift}(c)}{r} (B+cr) + \gamma^\alpha\frac{C_{frac}(c,\alpha)}{r^\alpha} (B+cr)^\alpha,
	\end{equation*}
	with $C_{drift}(c)>0$, $ C_{frac}(c,\alpha)>0$, for any $c\in (0,1)$.  
\end{proposition}

\begin{proof}
	By definition of the mollifier in \eqref{eq:mollifier} we have $0\leq \phi_r(x)\leq 1$ and $supp (\phi_r(x)) = \mathcal{B}_r(x^\star)$. Then 
	\begin{equation}\label{eq:mass_phi_relation}
		\rho_t(\mathcal{B}_r(x^\star)) \geq \int_{\RR^d} \phi_r(x) f(x,t) dx. 
	\end{equation}
	The idea is to find a lower bound for the right-hand side of \eqref{eq:mass_phi_relation}. Hence, we consider the weak form of \eqref{eq:fFP}
	\begin{equation*}
		\frac{d}{dt} \int_{\RR^d} \phi_r(x) f(x,t) dx = I_{drift}(t) + I_{frac}(t),
	\end{equation*}
	with
	\begin{equation*}
		\begin{split}
			I_{drift}(t) = & \nu \int_{\RR^d}  (\bar{x}_\mathcal{E}(t)-x) \nabla_x \phi_r(x) f(x,t) dx := \int_{\RR^d} T_1(x) f(x,t) dx \\
			I_{frac}(t) = & -\gamma^\alpha \int_{\RR^d} \vert \bar{x}_\mathcal{E}(t)-x\vert^\alpha (-\Delta)^{\alpha/2} \phi_r(x) f(x,t) dx := \int_{\RR^d} T_2(x) f(x,t) dx.
		\end{split}
	\end{equation*}
	Now we proceed by splitting the domain $\RR^d$ into different regions, depending on the sign of $T_1(x)$ and $T_2(x)$. 
	We begin by analyzing the drift term. We recall that 
	\begin{equation*}
		\nabla \phi_r(x) = \begin{cases}
			\frac{-2 r^2 (x-x^\star) }{(r^2 - \vert x-x^\star\vert^2)^2} \phi_r(x)  \qquad \text{if } \quad \vert x-x^\star \vert < r,\\
			0  \qquad  \qquad \qquad \qquad \qquad  \text{else}.
		\end{cases}
	\end{equation*}
	Hence, we get
	\begin{equation*}
		\text{sign}(T_1(x)) = - \text{sign}( \langle \bar{x}_\mathcal{E}(t)-x, x-x^\star\rangle). 
	\end{equation*}
	We define 
	\begin{equation*}
		\Omega_r = \{ x\in \RR^d : \vert x-x^\star\vert < r\},
	\end{equation*}
	and for any $c\in (0,1)$
	\begin{equation}
		\begin{split}
			& K_1 = \Omega_r^c \cup \{ x\in \Omega_r : \langle \bar{x}_\mathcal{E}(t)-x, x-x^\star\rangle \leq 0 \quad \text{and} \quad c r<\vert x-x^\star \vert < r \},\\
			& K_2 = \{ x\in \Omega_r : \langle \bar{x}_\mathcal{E}(t)-x, x-x^\star\rangle > 0 \quad \text{and} \quad \vert x-x^\star \vert \leq c r\}.
		\end{split}
	\end{equation}
	We note that $K_1 \cup K_2 = \RR^d$ and $K_1\cap K_2 = \emptyset$. 
	\paragraph{Case $K_1$.} 
	If $x\in \Omega_r$ and $\langle \bar{x}_\mathcal{E}(t)-x, x-x^\star\rangle\leq 0$ then 
	\begin{equation*}
		T_1(x) = \frac{-2 r^2 \nu}{(r^2 - \vert x-x^\star \vert^2)^2} \phi_r(x) \langle \bar{x}_\mathcal{E}(t)-x, x-x^\star\rangle \geq 0.
	\end{equation*}
	If $x \not \in \Omega_r$, then $\phi_r(x) = 0$ and $\nabla \phi_r(x)=0$, so $T_1(x)=0$. 
	Therefore, this yields
	\begin{equation*}
		\int_{K_1\cap \Omega_r} T_1(x) f(x,t) dx \geq 0. 
	\end{equation*}
	\paragraph{Case $K_2$.} If $x\in K_2$ then $\langle \bar{x}_\mathcal{E}(t)-x, x-x^\star\rangle> 0$ and $\vert x-x^\star \vert \leq cr$ for any $c\in (0,1)$.
	Now,
	\begin{equation*}
		\begin{split}
			T_1(x) = & -2 r^2 \nu \frac{\langle \bar{x}_\mathcal{E}(t)-x, x-x^\star\rangle}{(r^2-\vert x-x^\star\vert ^2)^2}\phi_r(x) \geq -2 r^2 \nu \frac{\vert \bar{x}_\mathcal{E}(t)-x\vert \vert x-x^\star\vert}{(r^2-\vert x-x^\star\vert ^2)^2} \phi_r(x) \geq \\ 
			& \geq -2 r^2 \nu \frac{\left( \vert \bar{x}_\mathcal{E}(t)-x^\star\vert + \vert x^\star - x \vert \right) \vert x-x^\star\vert}{(r^2-\vert x-x^\star\vert ^2)^2} \phi_r(x) \geq \frac{-2r^2\nu (B+cr) cr}{(r^2-c^2r^2)^2}\phi_r(x) = \\
			= & -\nu \frac{C_{drift}(c)}{r} (B+cr) \phi_r(x) : = -p_1 \phi_r(x),
		\end{split}
	\end{equation*}
	with $C_{drift}(c) = 2c/(1-c^2)^2$. Then, 
	\begin{equation}\label{eq:I_drift}
		I_{drift} = \int_{\RR^d} T_1(x) f(x,t) dx \geq \int_{K_2} T_1(x) f(x,t)dx \geq -p_1 \int_{\RR^d} \phi_r(x) f(x,t).
	\end{equation}
	We now focus on $I_{frac}$. 
	In analogy with the drift term, we split the domain into $\Omega_r^c$ and 
	\begin{align*}
		\Omega_{in} = \{x \in \Omega_r : \vert x-x^\star \vert \leq cr\}, \qquad
		\Omega_{bd} = \{x \in \Omega_r : cr< \vert x-x^\star \vert <r\}. 
	\end{align*}
	\paragraph{Case $\Omega_r^c$.} In $\Omega_r^c$, $\phi_r(x)=0$. Using the definition of fractional Laplacian 
	\begin{equation*}
		(-\Delta)^{\alpha/2} \phi_r(x) = c_\alpha PV \int_{\RR_d} \frac{\phi_r(x)-\phi_r(y)}{\vert x-y\vert^{1+\alpha}} dy = c_\alpha PV \int_{\RR_d} \frac{0-\phi_r(y)}{\vert x-y\vert^{1+\alpha}} dy \leq 0,
	\end{equation*}
	for $c_\alpha>0$
	\paragraph{Case $\Omega_{in}$.} 
	Define 
	\begin{equation}\label{eq:phi_1} 
		\phi_1(u) =   \begin{cases}
			\exp \left( 1- \frac{1}{1 - u^2} \right)  \qquad \text{if } \quad \vert u \vert < 1,\\
			0  \qquad  \qquad \qquad \qquad  ~ \text{else}. 
		\end{cases}  
	\end{equation}
	Note that $\phi_r(x) = \phi_1\left( \frac{x-x^\star}{r}\right)$. For any $x\in \RR^d$, set $u=\frac{x-x^\star}{r} $:
	\begin{equation}\label{eq:relation_phi1_phi_r} 
		\begin{split}
			(-\Delta)^{\alpha/2} \phi_r(x) =&   c_\alpha PV \int_{\RR_d} \frac{\phi_1(u)-\phi_1(\frac{y-x^\star}{r})}{\vert x-y\vert^{1+\alpha}} dy = \\ & = c_\alpha PV \int_{\RR_d} \frac{\phi_1(u)-\phi_1(z)}{(r\vert u-z\vert)^{1+\alpha}} r dz  = r^{-\alpha}  (-\Delta)^{\alpha/2} \phi_1(x).
		\end{split}
	\end{equation}
	If $x\in \Omega_{in}$ then $\vert x-x^\star\vert\leq cr$, which means $\vert ur\vert \leq cr$ that is $\vert u \vert \leq c$. Then on the compact set $[-c,c]^d \in (-1,1)^d$, $\phi_1\in \mathcal{C}_c^\infty$ and strictly positive. Hence it $\phi_1(u)$ attains a minimum $m(c)>0$, and $((-\Delta)^{\alpha/2}\phi_1(x))$  attains a maximum $M(\alpha,c)<+\infty$  respectively. 
	Then 
	\begin{equation}
		(-\Delta)^{\alpha/2} \phi_r(x) =  r^{-\alpha}  (-\Delta)^{\alpha/2} \phi_1(x) \leq  r^{-\alpha} \frac{M(\alpha,c)}{m(c)} m(c) \leq   r^{-\alpha} \frac{M(\alpha,c)}{m(c)}\phi_r(x), 
	\end{equation}
	where we recall that $\phi_1(u) = \phi_r(x)$ by definition.  
	\paragraph{Case $\Omega_{bd}$.} 
	Fix $d\in\mathbb{N}$ and $\alpha\in(1,2)$. Define the radial bump function $\phi_1(x)$ as in \eqref{eq:phi_1}. 	 
	Fix $\alpha\in(1,2)$ and consider $x$ such that $\eta<|x|<1$, for some $\eta\in(3/4,1)$ to be chosen later. Using the symmetric representation of the fractional Laplacian, we write
	\begin{equation*}
		(-\Delta)^{\alpha/2}\phi_1(x)
		= c_{d,\alpha}\int_{\mathbb{R}^d}\frac{2\phi_1(x)-\phi_1(x+h)-\phi_1(x-h)}{|h|^{d+\alpha}}\,dh,
	\end{equation*}
	so that it is enough to prove that the integral is negative.
	Define 
	\[
	m_0 := \min_{|y|\le 1/2}\phi_1(y)>0, 
	\qquad 
	\varepsilon_\eta := \sup_{\eta\le |x|<1}\phi_1(x)=\phi_1(\eta),
	\qquad 
	d_\eta := m_0-2\varepsilon_\eta.
	\]
	For $\eta$ sufficiently close to $1$, one has $d_\eta>0$.
	
	We split the integral into three regions. First, consider the translated set 
	\[ 
	S_x := x - \mathcal{B}_{1/2}(0) = \{ h = x-y : y\in \mathcal{B}_{1/2}(0)\}.
	\]   For $h\in S_x$ we have $\phi_1(x-h)\ge m_0$ and $\phi_1(x+h)=0$, hence
	\[
	2\phi_1(x)-\phi_1(x+h)-\phi_1(x-h) \le 2 \varepsilon_\eta-0-m_0 = -d_\eta<0.
	\]
	Moreover, $|h|\le 3/2$ on $S_x$, and since $\abs{S_x}=\omega_d 2^{-d+1}$, we conclude
	\begin{equation}\label{eq:contribution_1}
		\int_{S_x}\frac{2\phi_1(x)-\phi_1(x+h)-\phi_1(x-h)}{|h|^{d+\alpha}}\,dh
		\le - d_\eta\,\omega_d\,2^{-d+1}\Bigl(\frac{2}{3}\Bigr)^{d+\alpha}.
	\end{equation}
	Next, for $|h|<\rho$ with $\rho\le \eta-1/2$, a Taylor expansion yields
	\[
	|2\phi_1(x)-\phi_1(x+h)-\phi_1(x-h)| \le M_2 |h|^2,
	\]
	so that
	\begin{equation*}
		\left| \int_{|h|<\rho}\frac{2\phi_1(x)-\phi_1(x+h)-\phi_1(x-h)}{|h|^{d+\alpha}}\,dh\right| \leq M_2 \int_{|h|<\rho} \vert h \vert^{2-d-\alpha} dh  
		= M_2\,\omega_d\,\frac{\rho^{2-\alpha}}{2-\alpha}.
	\end{equation*}
	and in particular 
	\begin{equation}\label{eq:contribution_2}
		\int_{|h|<\rho}\frac{2\phi_1(x)-\phi_1(x+h)-\phi_1(x-h)}{|h|^{d+\alpha}}\,dh
		\le M_2\,\omega_d\,\frac{\rho^{2-\alpha}}{2-\alpha}.
	\end{equation}
	Finally, using $\phi_1\ge 0$, for $|h|\ge \rho$ we obtain
	\begin{equation}\label{eq:contribution_3}
		\int_{|h|\ge \rho}\frac{2\phi_1(x)-\phi_1(x+h)-\phi_1(x-h)}{|h|^{d+\alpha}}\,dh
		\le 2\varepsilon_\eta \int_{\abs{h}\geq \rho} \abs{h}^{-(d+\alpha)} dh = 2\varepsilon_\eta\,\omega_d\,\frac{\rho^{-\alpha}}{\alpha} .
	\end{equation}
	Combining the three contributions in \eqref{eq:contribution_1}-\eqref{eq:contribution_2}-\eqref{eq:contribution_3}, we get
	\begin{equation*}
		(-\Delta)^{\alpha/2}\phi_1(x)
		\le c_{d,\alpha}\,\omega_d\Bigg[
		- d_\eta\,2^{-d+1}\Bigl(\frac{2}{3}\Bigr)^{d+\alpha}
		+ \frac{M_2}{2-\alpha}\rho^{2-\alpha}
		+ \frac{2\varepsilon_\eta}{\alpha}\rho^{-\alpha}
		\Bigg].
	\end{equation*}
	Optimizing in $\rho>0$ gives the choice $\rho_*=\sqrt{2\varepsilon_\eta/M_2}$, for which
	\[
	\frac{M_2}{2-\alpha}\rho_*^{2-\alpha}
	+ \frac{2\varepsilon_\eta}{\alpha}\rho_*^{-\alpha}
	= \frac{2}{\alpha(2-\alpha)}(2\varepsilon_\eta)^{\frac{2-\alpha}{2}} M_2^{\frac{\alpha}{2}}.
	\]
	If $\rho_*\le \eta-1/2$, we deduce
	\begin{equation*}
		(-\Delta)^{\alpha/2}\phi_1(x)
		\le c_{d,\alpha}\,\omega_d\Bigg[
		- d_\eta\,2^{-d+1}\Bigl(\frac{2}{3}\Bigr)^{d+\alpha}
		+ \frac{2}{\alpha(2-\alpha)}(2\varepsilon_\eta)^{\frac{2-\alpha}{2}} M_2^{\frac{\alpha}{2}}
		\Bigg].
	\end{equation*}
	Thus, choosing $\eta$ sufficiently close to $1$ so that
	\begin{equation*}
		d_\eta\,2^{-d+1}\Bigl(\frac{2}{3}\Bigr)^{d+\alpha}
		>
		\frac{2}{\alpha(2-\alpha)}(2\varepsilon_\eta)^{\frac{2-\alpha}{2}} M_2^{\frac{\alpha}{2}},
	\end{equation*}
	we conclude that
	$(-\Delta)^{\alpha/2}\phi_1(x)<0$, for all  $\eta<|x|<1.
	$ 
	Now, since $(-\Delta)^{\alpha/2}\phi_r(x) = r^{-\alpha} (-\Delta)^{\alpha/2}\phi_1(x)$ as proved \eqref{eq:relation_phi1_phi_r}, we get that 
	\[
	(-\Delta)^{\alpha/2}\phi_r(x)<0 \qquad \text{for all } \eta<|x|<1.
	\]
	Therefore, gathering the previous estimates yields
	\begin{equation}\label{eq:I_frac}
		\begin{split}
			I_{frac} = &-\gamma^\alpha \int_{\RR^d} \vert \bar{x}_\mathcal{E} -x \vert^\alpha   (-\Delta)^{\alpha/2} \phi_r(x) f(x,t) dx \\
			& \geq -\gamma^\alpha \int_{\Omega_{in}} \vert \bar{x}_\mathcal{E} -x \vert^\alpha   (-\Delta)^{\alpha/2} \phi_r(x) f(x,t) dx \\
			& \geq -\gamma^\alpha (B+cr)^\alpha r^{-\alpha} C_{frac}(c,\alpha)  \int_{\RR^d} \phi_r(x) f(x,t) dx,% :=\\
			%&:= - p_2  \int_{\RR^d} \phi_r(x) f(x,t) dx,
		\end{split}
	\end{equation}
	with $C_{frac}(c,\alpha) =  M(\alpha,c)/m(c)$.
	Finally, by plugging together the drift and fractional diffusion parts \eqref{eq:I_drift}-\eqref{eq:I_frac} and by applying Gr{\"o}nwall inequality we get the desired mass bound in \eqref{eq:bound_mass}. 
\end{proof}

\end{document}
